% Structured FAQ (Generic Extra-Duty Assignment)
% XeLaTeX standards-compliant formatting

\documentclass[11pt,a4paper,oneside]{article}

% Language and encoding
\usepackage[english]{babel}

% Fonts for XeLaTeX
\usepackage{fontspec}
\setmainfont{Liberation Serif}
\setsansfont{Liberation Sans}
\setmonofont{Liberation Mono}

% Page layout
\usepackage[margin=1in]{geometry}
\usepackage{setspace}
\onehalfspacing

% Lists and formatting
\usepackage{enumitem}
\setlist[itemize]{noitemsep,topsep=0.5pt}
\setlist[enumerate]{noitemsep,topsep=0.5pt}

% Headers and footers
\usepackage{fancyhdr}
\pagestyle{fancy}
\fancyhf{}
\rfoot{\thepage}
\lfoot{FAQ: Extra-Duty Assignment (sample) ver. Sample}
\renewcommand{\footrulewidth}{0.4pt}
\renewcommand{\headrulewidth}{0pt}

% Custom plain style for title page: centered page number only, no rule
\fancypagestyle{plain}{
  \fancyhf{}
  \cfoot{\thepage}
  \renewcommand{\headrulewidth}{0pt}
  \renewcommand{\footrulewidth}{0pt}
}

% Document structure
\usepackage{titlesec}
\titleformat{\section}[block]{\Large\bfseries}{\thesection\quad}{0pt}{}
\titleformat{\subsection}[block]{\large\bfseries}{\thesubsection\quad}{0pt}{}
\titleformat{\subsubsection}[block]{\bfseries}{\thesubsubsection\quad}{0pt}{}

% Table of contents formatting
\usepackage{tocloft}
\renewcommand{\cftsecleader}{\cftdotfill{\cftdotsep}} % spaced-out periods
\renewcommand{\contentsname}{Table of Contents}

% Hyperlinks
\usepackage{hyperref}
\hypersetup{colorlinks=true, linkcolor=blue, urlcolor=blue, citecolor=blue}

% Graphics (logo) and TikZ (diagrams/flowcharts)
\usepackage{graphicx}
\usepackage{tikz}
\usetikzlibrary{shapes.geometric, arrows}

% Title page
\title{Extra-Duty Assignment FAQ\\Municipal Education Agency (Sample)}
\author{Human Resources Division}
\date{Effective Date: Sample}

\begin{document}
\maketitle

% Placeholder logo
\begin{center}
  \vspace{6pt}
  % \includegraphics[height=1.0in]{sample_logo.png}
  \textit{[Sample logo placeholder]}
\end{center}

\newpage
\tableofcontents
\newpage

% =====================
\section{FAQ Purpose}
\noindent Q1.1 \textbf{What is the purpose of this FAQ?}\\
\noindent A1.1 Lorem ipsum dolor sit amet, consectetur adipiscing elit. This sample FAQ provides generic guidance on extra-duty assignment policies, procedures, and requirements for staff of a municipal education agency. Sed do eiusmod tempor incididunt ut labore et dolore magna aliqua.\vspace{6pt}

% =====================
\section{General Topics}
\noindent Q2.1 \textbf{What is an extra-duty assignment?}\\
\noindent A2.1 Ut enim ad minim veniam, quis nostrud exercitation ullamco laboris. Extra-duty work is activity performed outside regular hours for which eligible employees are paid at an established hourly rate.\vspace{6pt}

\noindent Q2.2 \textbf{What is the assignment year?}\\
\noindent A2.2 Duis aute irure dolor in reprehenderit. The assignment year aligns with the agency fiscal year, from July 1 to June 30.\vspace{6pt}

\noindent Q2.3 \textbf{When can an employee perform extra-duty work?}\\
\noindent A2.3 Excepteur sint occaecat cupidatat non proident. For active employees, work occurs during non-regular hours. Retirees may work during standard hours if the posting allows.\vspace{6pt}

\noindent Q2.4 \textbf{Can a retiree work during the regular day?}\\
\noindent A2.4 Sunt in culpa qui officia deserunt mollit anim id est laborum. Yes, if specifically allowed by the advertisement and full-time staff have been considered first.\vspace{6pt}

% =====================
\section{Eligibility}
\noindent Q3.1 \textbf{Who can work extra-duty assignments?}\\
\noindent A3.1 Lorem ipsum dolor sit amet. Eligible pedagogical or instructional staff may participate. Certain support titles may be restricted.\vspace{6pt}

\noindent Q3.2 \textbf{Can substitute or temporary staff participate?}\\
\noindent A3.2 Consectetur adipiscing elit. Yes, after qualified full-time employees have been offered the opportunity.\vspace{6pt}

\noindent Q3.3 \textbf{Can a supervisor hold an extra-duty role?}\\
\noindent A3.3 Sed do eiusmod tempor. Yes, if the activity differs from regular duties and is performed outside regular hours.\vspace{6pt}

% =====================
\section{Advertisements and Postings}
\noindent Q4.1 \textbf{How are assignments advertised?}\\
\noindent A4.1 Ut labore et dolore magna aliqua. Postings occur at site, district, and central levels via the designated agency portal.\vspace{6pt}

\noindent Q4.2 \textbf{Must all activities be advertised?}\\
\noindent A4.2 Quis nostrud exercitation. Yes. Standard advertising periods apply unless an expedited process is approved with documented justification.\vspace{6pt}

% =====================
\section{Hours and Caps}
\noindent Q5.1 \textbf{What are typical annual hour limits?}\\
\noindent A5.1 Ullamco laboris nisi ut aliquip. Sample limits: 500 hours for supervisory titles, 400 hours for instructional staff. Waivers may extend these under review.\vspace{6pt}

\noindent Q5.2 \textbf{Are breaks required?}\\
\noindent A5.2 Ex ea commodo consequat. After five consecutive hours, an unpaid break is required. All time must be recorded.\vspace{6pt}

% =====================
\section{Contact}
\noindent Q6.1 \textbf{Whom do I contact for guidance?}\\
\noindent A6.1 Contact the sample Human Resources office at 555-0100 or staff@example.org. All content in this document is placeholder text for demonstration purposes only.\vspace{6pt}

\vspace{1cm}

\noindent\hrulefill

\noindent\textit{End of Sample Frequently Asked Questions Document}

\end{document}
